% concatenated from integral_cat.tex
\documentclass[a4paper,10pt,reqno]{amsart}

\usepackage{amsmath, amstext, paralist, amsthm, amssymb, epsfig, paralist, tikz, tikz-cd, graphicx}
\hoffset=20mm
\usepackage{hyperref}
\usepackage{xurl}

\newtheorem{thm}{Theorem}[section]
\newtheorem{lem}[thm]{Lemma}
\newtheorem{prop}[thm]{Proposition}
\newtheorem{ex}[thm]{Example}
\newtheorem{con}[thm]{Consequence}
\newtheorem{cor}[thm]{Corollary}


\newcommand{\myfootnote}[1]{
	\renewcommand{\thefootnote}{}
	\footnotetext{#1}
	\renewcommand{\thefootnote}{\arabic{footnote}}}

\renewcommand{\qedsymbol}{\ensuremath{\blacksquare}}
\renewcommand{\bar}{\overline}
\newcommand{\bc}[2]{\genfrac{[}{]}{0pt}{}{#1}{#2}}
\newcommand{\cb}[2]{{\scriptstyle[ #1\;#2]}}
\newcommand{\Ker}{\operatorname{Ker}\nolimits}
\newcommand{\Coker}{\operatorname{Cok}\nolimits}
\newcommand{\Cok}{\operatorname{Cok}\nolimits}
\newcommand{\cok}{\operatorname{cok}\nolimits}
\newcommand{\coker}{\operatorname{cok}\nolimits}
\newcommand{\Img}{\operatorname{Im}\nolimits}
\renewcommand{\Im}{\operatorname{Im}\nolimits}
\newcommand{\im}{\operatorname{im}\nolimits}
\newcommand{\coim}{\operatorname{coim}\nolimits}
\newcommand{\Coim}{\operatorname{Coim}\nolimits}
\newcommand{\id}{\operatorname{id}\nolimits}

\begin{document}
	\title{Some Remarks about Integral Categories}
	\author{Yaroslav Kopylov}
	\address{Yaroslav Kopylov\,
		%\orcidlink{0000-0002-0343-4424}\,
		%0000-0002-0343-4424
		\newline\hphantom{iii} Sobolev Institute of Mathematics
		\newline\hphantom{iii} 4 Koptyug Ave.
		\newline\hphantom{iii} 630090, Novosibirsk, Russia}
	\email{\href{mailto:yakop@math.nsc.ru}{yakop@math.nsc.ru}}
	
	\author{Max Zinchenko}
	\address{Max Zinchenko\,
		%\orcidlink{???}\,
		%???
		\newline\hphantom{iii} Novosibirsk State University
		\newline\hphantom{iii} 1 Pirogova St.
		\newline\hphantom{iii} 630090, Novosibirsk, Russia}
	\email{\href{mailto:m.zinchenko@g.nsu.ru}{m.zinchenko@g.nsu.ru}}
	
	\renewcommand{\thefootnote}{}
	\footnote{2010 \emph{Mathematics Subject Classification}: Primary 18A20; Secondary 46M18.}
	\footnote{\emph{Key words}: preabelian category, integral category, semi-abelian category, quasi-abelian category.}
	\vspace{-35pt}
	
	{\small
		\begin{abstract}\noindent{}
			We consider integral categories, a class of categories that  gained importance in the last three decades in the categorical foundations of algebra and functional analysis. We discuss some familiar criteria for integrality and prove new criteria for one- and two-sided integrality. For the first time in the literature, an example is given of a right but not left integral category.
	\end{abstract}}
	
	\maketitle
	
	\section{Introduction}\label{Introduction}\vspace{-5pt}
	
	Integral categories were first considered by B\u anic\u a and Popescu in 1965 in \cite{BanicaPopescu1965} under the~name {\em pr\'eab\'eliennes sp\'eciales}. These are preabelian categories (i.e., additive categories with kernels and cokernels) in which epimorphisms are stable under pullbacks and monomorphisms are stable under pushouts. In 2001, Rump called such categories integral and introduced the notions of a left integral category (a~preabelian one in which epimorphisms are stable under pullbacks) and a~right integral category (a~preabelian one in which monomorphisms are stable under pushouts).  
	
	Integral categories have recently been considered in representation theory (see, for example, \cite{BrHasTat2021, BuanMarsh2012, Nakaoka2013} and in functional analysis (see
	\cite{HSW2021, LawsonWegner2023}).
	
	Numerous examples and non-examples of integral categories can be found in \cite{HSW2021}. However, no example seems to have been known of a right integral category that is not left integral (or vice versa).  
	
	The structure of the paper is as follows. In Section~\ref{Prep}, 
	we give the necessary definitions and discuss the interrelationship between the types of preabelian categories we consider. In Section \ref{CrInt}, we discuss some familiar criteria for integrality and prove some extra criteria for one- and two-sided integrality of a~preabelian category. Then, in Section~\ref{ex-couple}, we adapt the results of Kopylov and Wegner \cite{KW2} on the derivation of exact couples in a semi-abelian category in the sense of Raikov (using the approach of Eckmann, Hilton \cite{EH} for constructing exact couples, see~\cite{M}), to integral categories. In the~last section, Section~\ref{Example}, we give an~explicit example of a~right integral category that is not left integral and discuss its consequences.
	
	\section{Notations and Preliminaries}\label{Prep}\vspace{-5pt}
	
	In the sequel, let $\mathcal{A}$ be a preabelian category, i.e.~an additive category with kernels and cokernels. For a morphism $f\colon E\rightarrow F$ in $\mathcal{A}$ we denote by $\ker f\colon \Ker f\rightarrow E$ its kernel and by $\cok f\colon F \rightarrow \Cok f$ its cokernel. Note that kernels and cokernels are unique only up to isomorphisms; we will however speak of \textit{the} kernel and \textit{the} cokernel of a given morphism $f$. We say that $f$ is a \textit{kernel} if there is a morphism $g$ such that $f=\ker g$. \textit{Cokernels} are defined dually. We denote by $\coim f\colon E\rightarrow \Coim f$ the cokernel of $\ker f$ and by $\im f\colon\Im f\rightarrow F$ the kernel of $\cok f$. As above, \textit{image} and \textit{coimage} are unique only up to isomorphims but we will also use definite articles in the sequel; the~same refers to pullbacks and pushouts. Any morphism $f$ in a~preabelian category admits a canonical decomposition $f=(\im f)\bar{f}\coim f$. Following B\u{a}nic\u{a}, Popescu \cite{BanicaPopescu1965} or Schneiders \cite[Definition 1.1.1]{Schneiders1999}, we say that $f$ is \textit{strict} if $\bar{f}$ is an isomorphism. By definition, $\mathcal{A}$ is abelian if and only if every morphism is strict.
	We say that a kernel $\alpha$ is semi-stable if all its pushouts along arbitrary morphisms are again kernels. Semi-stable cokernels are defined dually.
	
	\smallskip
	
	According to Rump \cite[p.~167]{Rump2001}, we say that $\mathcal{A}$ is \textit{left semi-abelian} if $\bar{f}$ is a monomorphism for each morphism $f$. Dually, we say that $\mathcal{A}$ is \textit{right semi-abelian} if $\bar{f}$ is an epimorphism for each morphism $f$. If $\mathcal{A}$ is left semi-abelian then each morphism $f$ admits a decomposition $f=ip$ with a cokernel $p$ and a monomorphism $i$. If dually $\mathcal{A}$ is right semi-abelian then each morphism $f$ admits a decomposition $f=ip$ with an epimorphism $p$ and a kernel $i$. In fact, the last assertions are even equivalent to the definitions of left and right semi-abelianity (see \cite[p.~167]{Rump2001}). A~left and right semi-abelian category
	is called a~semi-abelian category (in the~sense of Palamodov):~$\mathcal{A}$ is semi-abelian iff $\bar{f}$ is a \textit{bimorphism} (=\textit{regular morphism}), i.e.~$\bar{f}$ is a monomorphism and an epimorphism simultaneously, for each morphism $f$. 
	
	We need some results on preabelian categories. The first lemma summarizes well-known facts, see, e.g., \cite[Remark 1.1.2]{Schneiders1999}, \cite[Theorems 1 and 5]{RichmanWalker}, and \cite[Proposition 5.2]{Kelly1969}.
	
	\begin{lem}\label{l1} Let $\mathcal{A}$ be a preabelian category.
		\begin{itemize}
			\item[(i)] A morphism $f$ is a kernel if and only if $f=\im f$ and it is a cokernel if and only if $f=\coim f$.
			\item[(ii)] A morphism $f$ is strict if and only if there is a representation $f=f_1 f_0$ with a cokernel $f_0$ and a kernel $f_1$. In every such representation we have
			$f_0=\coim f$ and $f_1=\im f$.
			\item[(iii)]In every pullback
			$$
			\begin{tikzcd}
				\Ker p_G \arrow[r] & P \arrow[r, "p_G"] \arrow[d, "p_E"'] \arrow[r, "\text{PB}", phantom, bend right=60] & G \arrow[d, "t"]  &   \\
				\Ker f \arrow[r]   & E \arrow[r, "f"']                                                                   & F \arrow[r, "h"'] & H
			\end{tikzcd}
			$$
			we have $\ker f= p_E \ker p_G$. If $f$ is the kernel of a morphism $h$ then $p_G$ is the kernel of $ht$. If $f$ is a monomorphism then so is $p_G$.
			
			\smallskip
			
			Dually, in every pushout
			$$
			\begin{tikzcd}
				& E \arrow[r, "g"] \arrow[d, "s"'] & F \arrow[d, "s_F"] \\
				& G \arrow[r, "s_G"'] & S
				\arrow[from=1-2, to=2-3, phantom, "\text{\small PO}" description]
			\end{tikzcd}
			$$
			we have $\cok g=(\cok s_G) s_F$. If $g$ is the cokernel of a morphism $h$ then $s_G$ is the cokernel of $sh$. If $g$ is an epimorphism then so is $s_G$.\hfill\qed
		\end{itemize}
	\end{lem}
	
	We stress that the last statement of Lemma \ref{l1} implies that in any preabelian category, kernels pullback to kernels and cokernels pushout to cokernels.
	
	Following Rump \cite{Rump2001}, we say that a~preabelian category
	$\mathcal{A}$ is called {\em left integral} if pullbacks of epimorphisms in~$\mathcal{A}$ are epimorphisms. Dually,
	$\mathcal{A}$ is {\em right integral} if pushouts of monomorphisms are monomorphisms. 
	
	It is clear that a~left (resp., right) integral category is left (resp., right) semi-abelian. 
	
	A~left and right integral category is called {\em integral}. 
	
	A {\it left quasi-abelian category} is a preabelian category in which pullbacks of cokernels are cokernels; a {\em right quasi-abelian category} is a preabelian category in which pushouts of kernels are kernels. A left and right quasi-abelian category is called {\em quasi-abelian}.
	
	
	\section{Criteria for Integrality}\label{CrInt}
	
	In this section, we discuss some familiar criteria for integrality and prove some extra criteria for one- and two-sided integrality. 
	
	\smallskip
	
	Given a commutative square
	\begin{equation}\label{sq1}
		\begin{tikzcd}
			\Ker g \arrow[r] \arrow[d, "\hat{\alpha}"'] & C \arrow[r, "g"] \arrow[d, "\alpha"'] & D \arrow[d, "\beta"] \arrow[r] & \Cok g \arrow[d, "\hat{\beta}"] \\
			\Ker f \arrow[r]                            & A \arrow[r, "f"']                     & B \arrow[r]                    & \Cok f                         
		\end{tikzcd}
	\end{equation}
	in a preabelian category $\mathcal{A}$, we denote by $\hat\alpha\colon\Ker g\rightarrow\Ker f$ the unique morphism satisfying $\alpha\ker g=(\ker f)\hat\alpha$ and dually by $\hat\beta\colon\Coker g\rightarrow\Coker f$ the unique morphism satisfying $(\coker f)\beta=\hat\beta\coker g$. The above square will be used throughout the remainder of the section without repetition. 
	
	The following proposition gives us a~criterion for a~preabelian category to be integral, which will be used below:
	
	\begin{prop}\label{t1.1} 
		The following are equivalent for preabelian category $\mathcal{A}$:
		\begin{itemize}
			\item[(i)] $\mathcal{A}$ is integral.
			\item[(ii)] If ~\eqref{sq1} is a pushout then the canonical morphism $\hat\alpha: \Ker g \to \Ker f$ is an epimorphism, and if~\eqref{sq1}
			is a~pullback then the~canonical morphism $\hat\beta: \Coker g \to \Coker f$ is a~monomorphism.
		\end{itemize}
	\end{prop}
	
	\begin{proof}
		(i)$\Rightarrow$(ii) 
		Since $\mathcal A$ is integral, it is also semi-abelian. Thus, we can decompose $g$ as $g=g_1 g_0$ with $g_1$ a monomorphism and $g_0$ a cokernel. Consider the pushout
		$$
		\begin{tikzcd}
			C \arrow[r, "g_0"] \arrow[d, "\alpha"'] \arrow[r, "\text{PO}", phantom, bend right=60] & A' \arrow[d, "\beta'"] \arrow[r, "g_1", hook] & D \arrow[d, "\beta"] \\
			A \arrow[r, "f_0"'] \arrow[rr, "f"', bend right=49]                                    & B' \arrow[r, "f_1"', dashed]                  & B                   
		\end{tikzcd}
		$$
		We have $\beta g=\beta g_1 g_0=f\alpha$. Thus, there is a unique morphism $f_1\colon B'\rightarrow B$ with $f=f_1 f_0$ and $\beta g_1=f_1 \beta'$. It is easy to see (see, for instance, \cite[Lemma 5.1]{Kelly1969}) that
		$$
		\begin{tikzcd}
			A' \arrow[r, "g_1"] \arrow[d, "\beta' "'] & D \arrow[d, "\beta"] \\
			B' \arrow[r, "f_1"'] & B
			\arrow[from=1-1, to=2-2, phantom, "\scriptstyle\text{PO}" description]
		\end{tikzcd}
		$$
		is again a pushout and so $f_1$ is a monomorphism because $\mathcal{A}$ integral.\\
		Let $x\colon\Ker f\rightarrow X$ be a morphism such that $x\hat\alpha=0$. Consider the pushout
		$$
		\begin{tikzcd}
			\Ker g \arrow[r, "\hat{\alpha}"] \arrow[d] & \Ker f \arrow[r, "x"] \arrow[d, "\ker f"'] \arrow[r, "\text{PO}", phantom, bend right=60] & X \arrow[d, "z_2"] \\
			C \arrow[r, "\alpha"']                     & A \arrow[r, "z_1"']                                                                       & Z                 
		\end{tikzcd}
		$$
		We have $z_1\alpha\ker g=z_1(\ker f)\hat\alpha = z_2 x \hat\alpha=0$ and $g_0=\coker\ker g$. Thus, there exists a unique morphism $w\colon A'\rightarrow Z$ such that $z_1\alpha=w g_0$. Since $\beta' g_0=f_0\alpha$ is a pushout, there is a unique morphism $\sigma\colon B'\rightarrow Z$ such that $w=\sigma\beta'$ and $z_1=\sigma f_0$. But then $z_2 x=z_1 \ker f=\sigma f_0 \ker f=0$ because $0=f\ker f=f_1 f_0 \ker f$ and $f_1$ is a monomorphism. Since $\ker f$ is a monomorphism and the category is right integral, $z_2$ is a monomorphism, whence $x=0$ and $\hat\alpha$ is an epimorphism. 
		
		By duality, we see that $\hat\beta$ is a~monomorphism.
		\smallskip
		
		(ii)$\Rightarrow$(i) If (1) is a pushout with $g$ a monomorphism then $(\ker f)\hat\alpha = \alpha\ker g=0$. Since $\hat\alpha$ is an epimorphism, this gives $\ker f=0$, i.e.~$f$ is a monomorphism. Hence, $A$ is right integral. By duality, we conclude that $\mathcal{A}$ is also left integral. 
	\end{proof}
	
	The~following Proposition is well known (see~\cite[Proposition~7]{BanicaPopescu1965}, \cite[Proposition~6]{Rump2001}, and \cite[Proposition~4.1]{BuanMarsh2012}).
	
	\begin{prop}\label{sa-int1}
		Let $\mathcal{A}$ be a semi-abelian category. The~following are equivalent:
		\item[(i)] $\mathcal{A}$ is integral.
		\item[(ii)] Bimorphisms in $\mathcal{A}$ are stable under pullbacks.
		\item[(iii)] Bimorphisms in $\mathcal{A}$ are stable under pushouts.
		\item[(iv)] The~system of all bimorphisms in $\mathcal{A}$ admits
		a~calculus of left fractions and a~calculus of right fractions.
	\end{prop}
	
	In the~following proposition, we suggest a~criterion 
	for the~right (resp., left) integrality of a~right (resp., left) semi-abelian category. 
	
	\begin{prop}\label{sa-int2}
		A~right semi-abelian category $\mathcal{A}$ is right integral if and only if pushouts of monomorphisms along epimorphisms in~$\mathcal{A}$ are monomorphisms.
		
		Dually, a~left semi-abelian category $\mathcal{A}$ is left integral if and only if pullbacks of epimorphisms along monomorphisms in~$\mathcal{A}$ are epimorphisms.
	\end{prop}
	
	\begin{proof}
		We prove the~first assertion.
		
		If $\mathcal{A}$ is right integral then the~assertion is obvious.
		
		Suppose that pushouts of monomorphisms along epimorphisms in~$\mathcal{A}$ are monomorphisms and consider a~pushout
		$$
		\begin{tikzcd}
			C \arrow[r, "g"] \arrow[d, "\alpha"'] & D \arrow[d, "\beta"] \\
			A \arrow[r, "f"'] & B
			\arrow[from=1-1, to=2-2, phantom, "\scriptstyle\text{PO}" description]
		\end{tikzcd}
		$$
		with $g$ a~monomorphism. Since $\mathcal{A}$ is right semi-abelian,
		$\overline{\alpha}$ is an~epimorphism, and so $\alpha=\alpha_1 \alpha_0$, where $\alpha_1:A'\to A$ is a~kernel and $\alpha_0:C\to A'$ is an~epimorphism. Consider the~pushout
		$$
		\begin{tikzcd}
			C \arrow[r, "g"] \arrow[d, "\alpha_0"'] & D \arrow[d, "\beta_0"] \\
			A' \arrow[r, "f_0"'] & B'
			\arrow[from=1-1, to=2-2, phantom, "\scriptstyle\text{PO}" description]
		\end{tikzcd}
		$$
		Since $f\alpha_1\alpha_0 = f\alpha = \beta g$ and the~above square
		is a~pushout, there exists a~unique morphism $\beta_1:B'\to B$ such that $f\alpha_1 = \beta_1 f_0$ and $\beta=\beta_1\beta_0$. Since $\alpha_0$ is an~epimorphism, it is known and easy to check that
		$$
		\begin{tikzcd}
			A' \arrow[r, "f_0"] \arrow[d, "\alpha_1"'] & A \arrow[d, "\beta_1"] \\
			A \arrow[r, "f"'] & B
			\arrow[from=1-1, to=2-2, phantom, "\scriptstyle\text{PO}" description]
		\end{tikzcd}
		$$
		is also a~pushout. Since $\mathcal{A}$ is right semi-abelian and $\alpha_1$ is a~kernel, we conclude that the~last square is a~pullback (see \cite[Proposition~3.1]{KW}). Thus, $0=\Ker f_0\equiv \Ker f$, and $f$ is a~monomorphism. 
	\end{proof}
	
	
	\section{Exact couples in an~integral category}\label{ex-couple}
	
	An exact couple in $\mathcal{A}$ is a diagram of the form
	\begin{equation}\label{ex-coup}
		\begin{tikzcd}
			D \arrow[]{rr}{\alpha}& &D \arrow[]{ld}{\beta}\\
			&E \arrow[]{ul}{\gamma}
		\end{tikzcd}
	\end{equation}
	such that $\im\alpha=\ker\beta$, $\im\beta=\ker\gamma$ and $\im\gamma=\ker\alpha$. Using Lemma~1 in \cite{Yak1979}, it is not hard to see that the last equalities are equivalent to $\cok\alpha=\coim\beta$, $\cok\beta=\coim\gamma$ and $\cok\gamma=\coim\alpha$, respectively. For instance, $\coim\beta=\cok\im\alpha=\cok\alpha$ and the dual computation yield the first equivalence.
	
	\begin{thm}\label{THM-1}\cite[Theorem 1.]{KW2} Let $\mathcal{A}$ be semiabelian and consider the exact couple \eqref{ex-coup}. Assume that $\alpha$, $\beta$ and $\gamma$ are strict. Then the Eckmann-Hilton construction below gives rise to the following two diagrams
		\begin{equation}\label{der-coup}
			\begin{tikzcd}
				D_1 \arrow[]{rr}{\alpha_1}& &D_1 \arrow[]{ld}{\beta_1^-}\\
				&E_1^- \arrow[]{ul}{\gamma_1^-}
			\end{tikzcd}\hspace{20pt}
			%
			\begin{tikzcd}
				D_1 \arrow[]{rr}{\alpha_1}& &D_1 \arrow[]{ld}{\beta_1^+}\\
				&E_1^+ \arrow[]{ul}{\gamma_1^+}
			\end{tikzcd}
		\end{equation}
		which we call the \emph{left} resp.~the \emph{right derived couple}. The diagrams in \eqref{der-coup} have the following properties.\vspace{3pt}
		
		\begin{itemize}
			\item[(i)] Both diagrams are exact couples.
			
			\item[(ii)] With $\partial=\beta\gamma$ we get $H^-(E,\partial)=E_1^-$ and $H^+(E,\partial)=E_1^+$. Here, the \emph{left} resp.~\emph{right cohomology} is defined via $H^-(E,\partial)=\Cok(\theta\colon\Im\partial\rightarrow\Ker\partial)$ resp.~$H^+(E,\partial)=\Ker(\tau\colon\Cok\partial\rightarrow\Coim\partial)$, where $\theta$ and $\tau$ are the natural maps.
			
			\item[(iii)] There is a canonical bimorphism $\omega\colon E_1^-\rightarrow E_1^+$ satisfying the equations $\omega\beta_1^-=\beta_1^+$, $\gamma_1^-=\gamma_1^+\omega$ and $(\ker\tau)\omega(\cok\theta)=(\cok\partial)(\ker\partial)$. The morphism $\omega$ is uniquely determined by the third equation and it is an isomorphism if $\ker\partial$ or $\cok\partial$ is semistable.
		\end{itemize}
	\end{thm}
	
	\begin{thm}\label{THM-2}\cite[Theorem 2.]{KW2} Let $\mathcal{A}$ be semiabelian and consider the exact couple \eqref{ex-coup}. Assume that $\beta$ and $\gamma$ are strict and that $\ker\gamma$ and $\cok\beta$ are semistable. If the powers $\alpha^k$ are strict for $1\leqslant{}k\leqslant{}n$, then the derivation process of Theorem \ref{THM-1} can be performed $n$ times, i.e., the Eckmann-Hilton construction gives rise to a complete full binary tree of depth $n$ consisting of exact couples.
	\end{thm}
	
	Here we specify these theorems for integral categories. Namely, the~following assertions hold: 
	
	\begin{thm}\label{COR-1}
		If, in Theorem \ref{THM-1}, the category $\mathcal{A}$ is assumed to be integral, then the conclusion remains valid for arbitrary morphisms $\gamma$ and $\beta$, not necessarily strict.
	\end{thm}
	
	\begin{thm}\label{COR-2}
		If, in Corollary \ref{COR-1}, the category $\mathcal{A}$ is assumed to be quasi-abelian, then the right and left derived couples coincide.
	\end{thm}
	
	\smallskip
	
	Let $\mathcal{A}$ be preabelian. We fix the exact couple \eqref{ex-coup} and assume that the morphism $\alpha$ is strict, i.e., it has a decomposition $\alpha=\rho\sigma$ with a kernel $\rho$ and a cokernel $\sigma$. Taking the exactness into account it follows that $\rho=\ker\beta$ and $\sigma=\cok\gamma$ and we can consider the diagram
	\begin{equation}\label{eck-hil1}
		\begin{tikzcd}
			D_1  & &D_1 \arrow[tail]{d}{\rho}\\
			D \arrow[two heads]{u}{\sigma} \arrow[]{r}{\beta} & E \arrow[]{r}{\gamma}&  D
		\end{tikzcd}
	\end{equation}
	with $D_1=\Im\alpha=\Coim\alpha=\Ker\beta=\Cok\gamma$. From \eqref{eck-hil1} we get the two diagrams
	\begin{equation}\label{eck-hil2}
		\begin{tikzcd}[column sep=1.1ex, row sep=1.1ex]
			D_1\arrow[]{rr}{\beta_1^-}  && E_1^-\arrow[]{rr}{\gamma_1^-} &&D_1 \arrow[equal]{dd}\\
			&\text{\footnotesize\rm PO}&&\text{\footnotesize\rm\phantom{PB}}&\\
			D\arrow[two heads]{uu}{\sigma}\arrow[]{rr}{\beta'}   && E_{\rho}\arrow[two heads]{uu}[swap]{\sigma'}\arrow[]{rr}{\gamma'}\arrow[tail]{dd}{\rho'} &&D_1 \arrow[tail]{dd}{\rho}\\
			&\text{\footnotesize\rm\phantom{PB}}&&\text{\footnotesize\rm PB}&\\
			D \arrow[equal]{uu} \arrow[]{rr}{\beta} && E \arrow[]{rr}{\gamma}&&  D
		\end{tikzcd}\hspace{30pt}
		%
		\begin{tikzcd}[column sep=1.1ex, row sep=1.1ex]
			D_1\arrow[]{rr}{\beta_1^+} && E_1^+\arrow[]{rr}{\gamma_1^+}\arrow[tail]{dd}{\rho''}&& D_1 \arrow[tail]{dd}{\rho}\\
			&\text{\footnotesize\rm\phantom{PB}}&&\text{\footnotesize\rm PB}&\\
			D_1\arrow[equal]{uu}\arrow[]{rr}{\beta''} && E^{\sigma}\arrow[]{rr}{\gamma''} &&D\arrow[equal]{dd} \\
			&\text{\footnotesize\rm PO}&&\text{\footnotesize\rm\phantom{PB}}&\\
			D \arrow[two heads]{uu}{\sigma} \arrow[]{rr}{\beta} && E \arrow[]{rr}{\gamma}\arrow[two heads]{uu}[swap]{\sigma''}&&  D
		\end{tikzcd}
	\end{equation}
	by the classical construction of Eckmann, Hilton \cite{EH}. We start with the left diagram. We first form the pullback to obtain $\gamma'$ and $\rho'$. Then we apply its universal property and get $\beta'$ with $\gamma'\beta'=0$ and $\rho'\beta'=\beta$. Next, we form the pushout to obtain $\sigma'$ and $\beta_1^-$. By the universal property of the pushout we finally get $\gamma_1^-$ with $\gamma_1^-\beta_1^-=0$ and $\gamma_1^-\sigma'=\gamma'$. By \cite[Lemma~2.1(iii)]{KW}, $\rho'$ is a kernel and $\sigma'$ is a cokernel. The right diagram is constructed dually. The diagrams in \eqref{der-coup} are obtained by setting $\alpha_1=\sigma\rho$.
	
	\smallskip
	
	\begin{lem}\label{second}\cite[Lemmas 3 and 4]{KW2} Let $\mathcal{A}$ be left semiabelian and let $\alpha$, $\beta$ and $\rho$ be morphisms. If $\alpha=\cok\beta$ and $\im\beta=\ker(\rho\alpha)$ then $\rho$ is a monomorphism.
		
		Dually, let $\mathcal{A}$ be right semiabelian and let $\alpha$, $\beta$ and $\rho$ be morphisms. If $\coim\alpha=\cok(\rho\beta)$ and $\rho$ is a kernel then $\im\beta=\ker(\alpha\rho)$.
	\end{lem}
	
	Now, we prove Theorems \ref{COR-1} and \ref{COR-2}.
	
	\begin{proof}\textit{(of Theorem \ref{COR-1})} (i) We show $\im\beta_1^-=\ker\gamma_1^-$, $\im\alpha_1=\ker\beta_1^-$ and $\im\gamma_1^-=\ker\alpha_1$. By Lemma~\ref{second}, $\cok\beta=\coim\gamma$ and $\beta=\rho'\beta'$ imply that $\im\beta'=\ker(\gamma\rho')$. Hence, $\im\beta'=\ker(\rho\gamma')=\ker\gamma'$ and therefore $\coim\gamma'=\cok\im\beta'=\cok\beta'$. By \cite[Lemma 2.1(iii)]{KW}, we have $\cok\beta'=(\cok\beta_1^-)\sigma'$. In addition $\gamma'=\gamma_1^-\sigma'$ and $\sigma'$ is a cokernel. Therefore, $\coim\gamma'=(\coim\gamma_1^-)\sigma'$ by \cite[Cor.~2.3(ii) and Proposition~3.2]{KW}. We compute $(\cok\beta_1^-)\sigma'=\cok\beta'=\coim\gamma'= (\coim\gamma_1^-)\sigma'$ which gives $\cok\beta_1^-=\coim\gamma_1^-$ since $\sigma'$ is epic. This is equivalent to $\im\beta_1^-=\ker\gamma_1^-$.
		\smallskip
		\\We have $\beta=\rho'\beta'$, $\beta$ is strict and $\rho'$ is a kernel. Thus, $\rho'\beta'=\im(\rho'\beta')\bar{\beta}\coim(\rho'\beta')=\rho'(\im\beta')\bar{\beta}\coim(\rho'\beta')$ where the last equality follows from \cite[Corollary~2.3(i) and Proposition~3.1]{KW} and $\bar{\beta}$ is an isomorphism. We get $\beta'=(\im\beta')\bar{\beta}\coim(\rho'\beta')$ and thus $\beta'$ is strict by \cite[Remark~1.1.2(c)]{Schneiders1999}. Now we can use Proposition \ref{t1.1} to conclude that the morphism $\widehat\sigma$ defined by the equality $\sigma(\ker\beta')=(\ker\beta_1^-)\widehat\sigma$ is an epimorphism. We have $\ker\beta'=\ker(\rho'\beta')=\ker\beta=\rho$ and hence $\alpha_1=\sigma\rho=(\ker\beta_1^-)\widehat\sigma$. Therefore, $\im\alpha_1=\ker\beta_1^-$ by \cite[Corollary~2.3(i) and Proposition~3.1]{KW}.
		
		\smallskip
		
		By dualizing the last paragraph, we obtain $\cok\gamma_1^+=\coim\alpha_1$ and thus $\cok\gamma_1^-=\cok(\gamma_1^+\omega)=\cok\gamma_1^+$ with $\omega$ as in Theorem~\ref{THM-1}. Therefore, $\cok\gamma_1^-=\coim\alpha_1$ which is equivalent to $\im\gamma_1^-=\ker\alpha_1$.
		
		\smallskip
		
		The exactness of the second couple can be proved dually. The rest of the proof is exactly the same as that of \cite[Theorem 1.]{KW2}.
	\end{proof}
	
	\begin{proof}\textit{(of Theorem \ref{COR-2})}
		Since in a quasi-abelian category kernels are stable under pushouts and cokernels are stable under pullbacks, both $\ker \partial$ and $\coker \partial$ are, in particular, semistable. Hence, by \cite[Theorem~1(iii)]{KW2}, the associated derived couples coincide.
	\end{proof}
	
	
	\section{An Example of one-sided integrality}\label{Example}
	
	In this section, we present an elegant example of a non-semi-abelian preabelian category
	which was considered by Jeremy Rickard on Math StackExchange.${}^1$\footnote{%
		${}^1$ \url{https://math.stackexchange.com/questions/3631067/example-of-a-pre-abelian-category-but-not-a-semi-abelian-category}}
	Below we prove that it is right integral but not left integral.
	Since a~semi-abelian category is left integral if and only if it is right integral \cite[p.~12, Corollary]{Rump2001}, such a~category is necessarily not left semi-abelian.
	
	Let $\mathcal{C}$ be the category of diagrams
	$$
	U \xrightarrow{\alpha} V \xrightarrow{\beta} W
	$$
	of vector spaces over a fixed field $k$. Since the category of vector spaces over a fixed field is abelian, it follows that $\mathcal{C}$ is preabelian because kernels and cokernels are computed componentwise:
	$$
	\begin{tikzcd}
		\Ker f \arrow[d] \arrow[r, "\alpha''"] & \Ker g \arrow[d] \arrow[r, "\beta''"] & \Ker h \arrow[d] \\
		U \arrow[d, "f"] \arrow[r, "\alpha"]   & V \arrow[d, "g"] \arrow[r, "\beta"]   & W \arrow[d, "h"] \\
		U' \arrow[r, "\alpha'"] \arrow[d]      & V' \arrow[r, "\beta'"] \arrow[d]      & W' \arrow[d]     \\
		\Coker f \arrow[r]                     & \Coker g \arrow[r]                    & \Coker h.
	\end{tikzcd} \eqno(1)
	$$
	A morphism $(f,g,h)$ in $\mathcal{C}$ is strict if and only if each of the morphisms $f$, $g$, and $h$ is strict in the category of vector spaces. Hence $\mathcal{C}$ is an abelian category.
	
	Let $\mathcal{D}$ be the full subcategory of $\mathcal{C}$ consisting of those diagrams for which
	$$
	V=\im \alpha+\ker \beta.
	$$
	The subcategory $\mathcal{D}$ is closed under quotients in $\mathcal{C}$, and therefore cokernels in $\mathcal{D}$ exist and coincide with those in $\mathcal{C}$. It follows that $\mathcal{D}$ also has kernels. Indeed, let
	$$
	\Ker f \xrightarrow{\alpha''} \Ker g \xrightarrow{\beta''} \Ker h
	$$
	be the kernel of a morphism $(f,g,h)$ in $\mathcal{C}$. Then the kernel in $\mathcal{D}$ is given by
	$$
	\Ker f \xrightarrow{\alpha''} \im \alpha''+\ker \beta'' \xrightarrow{\beta''} \Ker h.
	$$
	Thus, $\mathcal{D}$ is a preabelian category.
	
	Now consider the following nonzero morphism:
	$$
	\begin{tikzcd}
		k \arrow[r, equal] \arrow[d, equal] & k \arrow[r, equal] \arrow[d] & k \arrow[d] \\
		k \arrow[r]                         & 0 \arrow[r]                  & 0.
	\end{tikzcd}
	$$
	Its coimage in $\mathcal{D}$ is
	$$k \to k \to 0,$$
	whereas its image is
	$$k \to 0 \to 0.$$
	Therefore, the canonical morphism from the coimage to the image is not a monomorphism. Hence $\mathcal{D}$ is not left semi-abelian, and therefore it is left integral.
	\par Observe that a morphism $(f,g,h)$ is a monomorphism in $\mathcal{D}$ if and only if it is a monomorphism in $\mathcal{C}$.
	
	Indeed, suppose that $(f,g,h)$ is a monomorphism in $\mathcal{D}$, and consider the diagram $(1)$ where
	$$
	\Ker f \xrightarrow{\alpha''} \Ker g \xrightarrow{\beta''} \Ker h
	$$
	is the kernel of $(f,g,h)$ in $\mathcal{C}$. Then
	$$
	\Ker f \xrightarrow{\alpha''} \bigl(\im \alpha''+\ker \beta''\bigr) \xrightarrow{\beta''} \Ker h
	$$
	is the kernel of $(f,g,h)$ in $\mathcal{D}$. Since $(f,g,h)$ is a monomorphism in $\mathcal{D}$, it follows that
	$$
	\Ker f=\im \alpha''+\ker \beta''=\Ker h=0.
	$$
	Since $\Ker f=0$, we have $\im \alpha''=0$, and therefore $\ker \beta''=0$. If $\Ker g\neq 0$, then
	$$
	\ker\bigl(\Ker g \xrightarrow{\beta''} 0\bigr)=\Ker g\neq 0,
	$$
	which is impossible. Hence $\Ker g=0$. Therefore $(f,g,h)$ is a monomorphism in $\mathcal{C}$. 
	\par We next show that $\mathcal{D}$ is closed under direct sums.\\
	Let
	$$
	X=(U_1 \xrightarrow{\alpha_1} V_1 \xrightarrow{\beta_1} W_1),
	\qquad
	Y=(U_2 \xrightarrow{\alpha_2} V_2 \xrightarrow{\beta_2} W_2)
	$$
	be objects of $\mathcal{D}$. Then
	$$
	X\oplus Y=
	\bigl(
	U_1\oplus U_2 \xrightarrow{\alpha_1\oplus \alpha_2}
	V_1\oplus V_2 \xrightarrow{\beta_1\oplus \beta_2}
	W_1\oplus W_2
	\bigr).
	$$
	Moreover,
	$$
	\im(\alpha_1\oplus \alpha_2)+\ker(\beta_1\oplus \beta_2)
	=
	(\im \alpha_1\oplus \im \alpha_2)+(\ker \beta_1\oplus \ker \beta_2),
	$$
	and therefore
	$$
	\im(\alpha_1\oplus \alpha_2)+\ker(\beta_1\oplus \beta_2)
	=
	(\im \alpha_1+\ker \beta_1)\oplus (\im \alpha_2+\ker \beta_2).
	$$
	Since $X,Y\in\mathcal D$, we have
	$$
	\im \alpha_1+\ker \beta_1=V_1,
	\qquad
	\im \alpha_2+\ker \beta_2=V_2.
	$$
	Hence
	$$
	\im(\alpha_1\oplus \alpha_2)+\ker(\beta_1\oplus \beta_2)=V_1\oplus V_2,
	$$
	so $X\oplus Y\in\mathcal D$. Thus $\mathcal D$ is closed under direct sums.
	
	Consequently, pushouts of morphisms in $\mathcal D$ may be computed in $\mathcal C$. Since monomorphisms are preserved and reflected, it follows that the pushout of any monomorphism is again a monomorphism. So $\mathcal{D}$ is not left integral but right integral.
	
	We now show that $\mathcal{D}$ is right quasi-abelian.
	
	Let $i\colon U\to X$ be a kernel in $\mathcal{D}$, and consider its pushout along an arbitrary morphism $g\colon U\to Z:$
	$$
	\begin{tikzcd}
		U \arrow[r, "i"] \arrow[d, "g"] \arrow[r, "\mathrm{PO}", phantom, bend right=60] & X \arrow[d] \\
		Z \arrow[r, "j"] & Y.
	\end{tikzcd}
	$$
	Since $\mathcal{C}$ is abelian, it follows that $j$ is a kernel in $\mathcal{C}$. Hence $\im_{\mathcal{C}} j = j,$ and
	$$
	\Coker_{\mathcal{C}} j=\Coker_{\mathcal{D}} j,
	$$
	because $\mathcal{D}$ is closed under quotients.
	
	Consider the canonical decomposition of $j$ in $\mathcal{C}$:
	$$
	\begin{tikzcd}
		Z \arrow[r, "j"] \arrow[rd, "\cong"'] & Y \arrow[r] & \Coker_{\mathcal{C}}j \\
		& \Im_{\mathcal{C}}j \arrow[u] &
	\end{tikzcd}
	$$
	Since both $\Im_{\mathcal{C}} j$ and $Y$ belong to $\mathcal{D}$, and $\mathcal{D}$ is a full subcategory of $\mathcal{C}$, we conclude that
	$$
	\im_{\mathcal{D}} j=\im_{\mathcal{C}} j=j.
	$$
	Therefore, $j$ is a kernel in $\mathcal{D}$. This proves that $\mathcal{D}$ is right quasi-abelian.
	
	Using this example, we obtain an example of a category that is right integral, but neither right quasi-abelian nor left semi-abelian. Indeed, it suffices to consider the product category $\mathcal{D}\times \mathbf{BOR}$ since, for product categories, such properties are inherited componentwise \cite[Theorem 3.8]{HSW2021}.
	
	
	\section*{Acknowledgements}
	The work of Ya.~Kopylov was carried out in the framework of the State Task to the
	Sobolev Institute of Mathematics (Project FWNF--2026--0026). 
	
	We are very grateful to
	Professor Sven-Ake Wegner for carefully reading the manuscript and
	pointing out several remarks.
	
	%The authors would like to thank the referee for his careful work.
	
	
	\begin{thebibliography}{100}
		
		\bibitem{BanicaPopescu1965}
		C.~B\u anic\u a and N.~Popescu, \emph{Sur les cat\'egories pr\'eab\'eliennes},
		Rev. Roumaine Math. Pures Appl. \textbf{10} (1965), 621--633.
		
		\bibitem{BrHasTat2021}
		T. Br\"{u}stle, S. Hassoun, and A. Tattar, \emph{Intersections, sums, and the Jordan-H¨older property for exact categories}, J.
		Pure Appl. Algebra \textbf{225} (2021), no. 11, 106724.
		
		\bibitem{BuanMarsh2012}
		A.~B.~Buan, B.~R.~Marsh, 
		\emph{From triangulated categories to module categories via localisation. II: Calculus of fractions}, 
		J. Lond. Math. Soc., II. Ser. \textbf{86}, 152--170 (2012); corrigendum \textbf{87}, 643 (2013).
		
		\bibitem{EH}
		B.~Eckmann and P.~J. Hilton, \emph{Exact couples in an abelian category}, J.
		Algebra \textbf{3} (1966), 38--87.
		
		\bibitem{HSW2021}
		S. Hassoun, A. Shah and S.-A. Wegner (2021)
		\emph{Examples and non-examples of integral categories and the admissible intersection property},
		Cahiers Topologie Géom. Différentielle Catég. 62, no. 3, 329-354
		
		\bibitem{Kelly1969}
		G.~M. Kelly, \emph{Monomorphisms, epimorphisms, and pull-backs}, J. Austral.
		Math. Soc. \textbf{9} (1969), 124--142.
		
		\bibitem{Kopylov2004}
		Ya. Kopylov, \emph{Exact couples in a {R}a\u{\i}kov semi-abelian category},
		Cah. Topol. G\'eom. Diff\'er. Cat\'eg. \textbf{45} (2004), no.~3, 162--178.
		
		\bibitem{KW}
		Ya. Kopylov and S.-A. Wegner, \emph{On the notion of a semi-abelian category in
			the sense of {P}alamodov}, Appl. Categor. Struct. \textbf{20} (2012), no.~5,
		531--541.
		
		\bibitem{KW2}
		Ya. Kopylov and S.-A. Wegner, \emph{Exact couples in semiabelian categories revisited}, J. Algebra \textbf{414} (2014), 264--270. doi{10.1016/j.jalgebra.2014.04.030}
		
		\bibitem{LawsonWegner2023}
		M. Lawson and S.~A. Wegner, 
		\emph{The category of Silva spaces is not integral}, 
		Homology Homotopy Appl. \textbf{25}, 367--374 (2023).
		
		\bibitem{M}
		W.~S. Massey, \emph{Exact couples in algebraic topology. {I}, {II}}, Ann. of
		Math. (2) \textbf{56} (1952), 363--396.
		
		\bibitem{Nakaoka2013}
		H. Nakaoka, \emph{General heart construction for twin torsion pairs on triangulated categories}, J. Algebra \textbf{374} (2013),
		195--215.
		
		\bibitem{RichmanWalker}
		F.~Richman and E.~A. Walker, \emph{Ext in pre-{A}belian categories}, Pacific J.
		Math. \textbf{71} (1977), no.~2, 521--535.
		
		\bibitem{Rump2001}
		W.~Rump, \emph{Almost abelian categories}, Cahiers Topologie G\'eom.
		Diff\'erentielle Cat\'eg. \textbf{42} (2001), no.~3, 163--225.
		
		\bibitem{Rump2008}
		W.~Rump \emph{A counterexample to {R}aikov's conjecture}, Bull. Lond. Math.
		Soc. \textbf{40} (2008), no.~6, 985--994.
		
		\bibitem{Schneiders1999}
		J.-P. Schneiders, \emph{Quasi-{A}belian categories and sheaves}, M\'{e}m. Soc.
		Math. Fr., Nouv. S\'{e}r. \textbf{76}, 1999.
		
		\bibitem{SiegWegner2011}
		D.~Sieg and S.-A. Wegner, \emph{Maximal exact structures on additive
			categories}, Math. Nachr. \textbf{284} (2011), no.~16, 2093--2100.
		
		\bibitem{Yak1979}
		A.~V. Yakovlev, \emph{Homological algebra in pre-abelian categories}, Zap.
		Nauchn. Sem. Leningrad. Otdel. Mat. Inst. Steklov. (LOMI) \textbf{94} (1979),
		131--141, English translation: J. Soviet Math., \textbf{19} (1982), 1060--1067.
		
		
	\end{thebibliography}
	
\end{document}



{\renewcommand{\baselinestretch}{1}\normalsize\small
	%\bibliography{BIB-2}}
\def\cprime{$'$}
\providecommand{\bysame}{\leavevmode\hbox to3em{\hrulefill}\thinspace}
\providecommand{\MR}{\relax\ifhmode\unskip\space\fi MR }
% \MRhref is called by the amsart/book/proc definition of \MR.
\providecommand{\MRhref}[2]{%
	\href{http://www.ams.org/mathscinet-getitem?mr=#1}{#2}
}
\providecommand{\href}[2]{#2}
\begin{thebibliography}{10}\vspace{-10pt}
	
	\bibitem{Grandis2010}
	M.~Grandis, \emph{Homotopy spectral sequences}, J. Homotopy Relat. Struct.
	\textbf{5} (2010), no.~1, 213--252.
	
	\bibitem{BanicaPopescu1965}
	C.~B\u anic\u a and N.~Popescu, \emph{Sur les cat\'egories pr\'eab\'eliennes},
	Rev. Roumaine Math. Pures Appl. \textbf{10} (1965), 621--633.
	\vspace{-7pt}
	\bibitem{BondalVanDenBergh}
	A.~Bondal and M.~van~den Bergh, \emph{Generators and representability of
		functors in commutative and noncommutative geometry}, Mosc. Math. J.
	\textbf{3} (2003), no.~1, 1--36, 258.
	\vspace{-7pt}
	\bibitem{BonetDierolf2005}
	J.~Bonet and S.~Dierolf, \emph{The pullback for bornological and
		ultrabornological spaces}, Note Mat. \textbf{25} (2005/06), no.~1, 63--67.
	\vspace{-7pt}
	\bibitem{Buehler2008}
	T.~B\"uhler, \emph{On the algebraic foundations of bounded cohomology}, Mem. Amer. Math. Soc., DOI: 10.1090/S0065-9266-2011-00618-0
	\vspace{-7pt}
	\bibitem{Buehler2009}
	\bysame, \emph{Exact categories}, Expo. Math. \textbf{28} (2010), no.~1, 1--69.
	\vspace{-7pt}
	\bibitem{BD}
	B.~Dierolf, \emph{Exaktheitsstrukturen}, Manuscript, University of Trier, 2010.
	\vspace{-7pt}
	\bibitem{Floret}
	K.~Floret, \emph{Some aspects of the theory of locally convex inductive limits}, in: Functional Analysis: Surveys and Recent Results II, K.D.~Bierstedt and B.~Fuchssteiner (eds.), North-Holland Math.~Stud. \textbf{38}, North-Holland, Amsterdam, 1980, 205--237.
	\vspace{-7pt}
	\bibitem{FloretWloka}
	K.~Floret and J.~Wloka, \emph{Einf\"uhrung in die {T}heorie der lokalkonvexen
		{R}\"aume}, Lecture Notes in Mathematics, No. 56, Springer-Verlag, Berlin,
	1968.
	\vspace{-7pt}
	\bibitem{Grand91-1}
	M.~Grandis, \emph{On homological algebra, {I}.~Semiexact and homological
		categories}, Dip. Mat. Univ. Genova, Preprint, 1991, available at http://www.dima.unige.it/~grandis/\linebreak{}HA1.Pr1991(Rev.09).ps
	\vspace{-7pt}
	\bibitem{Grothendieck1954}
	A.~Grothendieck, \emph{Sur les espaces ($\mathcal{F}$) et ($\mathcal{DF}$)},
	Summa Brasil. Math. \textbf{3} (1954), 57--123.
	\vspace{-7pt}
	\bibitem{JanelidzeMarkiTholen2002}
	G.~Janelidze, L.~M{\'a}rki, and W.~Tholen, \emph{Semi-abelian categories}, J.
	Pure Appl. Algebra \textbf{168} (2002), no.~2--3, 367--386.
	\vspace{-7pt}
	\bibitem{Jarchow}
	H.~Jarchow, \emph{Locally {C}onvex {S}paces}, B. G. Teubner, Stuttgart, 1981.
	\vspace{-7pt}
	\bibitem{Kashiwara}
	M.~Kashiwara, \emph{Equivariant derived category and representation of real
		semisimple {L}ie groups}, Representation theory and complex analysis, Lecture
	Notes in Math., vol. 1931, Springer, Berlin, 2008, pp.~137--234.
	\vspace{-7pt}
	\bibitem{Kelly1969}
	G.~M. Kelly, \emph{Monomorphisms, epimorphisms, and pull-backs}, J. Austral.
	Math. Soc. \textbf{9} (1969), 124--142.
	\vspace{-7pt}
	\bibitem{Kopylov2009}
	Ya.~A. Kopylov, \emph{The five- and nine-lemmas in {P}-semi-abelian
		categories}, Sibirsk. Mat. Zh. \textbf{50} (2009), no.~5, 1097--1104, English
	translation: Siberian Math. J. {\bf 50} (2009), no.~5, 867--873.
	\vspace{-7pt}
	\bibitem{KoK00}
	Ya.~A. Kopylov and V.~I. Kuz{\cprime}minov, \emph{Exactness of the cohomology
		sequence for a short exact sequence of complexes in a semiabelian category},
	Siberian Adv. Math. \textbf{13} (2003), no.~3, 72--80.
	\vspace{-7pt}
	\bibitem{KoK09}
	\bysame, \emph{The {K}er-{C}oker sequence and its generalization in some
		classes of additive categories}, Sibirsk. Mat. Zh. \textbf{50} (2009), no.~1,
	107--117, English translation: Siberian Math. J. \textbf{50} (2009), no.~1,
	86--95.
	\vspace{-7pt}
	\bibitem{KCh}
	V.~I. Kuz{\cprime}minov and A.~Yu. Cherevikin, \emph{Semiabelian categories},
	Sibirsk. Mat. Zh. \textbf{13} (1972), 1284--1294, 1420, English translation:
	Siberian Math. J. \textbf{13} (1972), no.~6, 895--902.
	\vspace{-7pt}
	\bibitem{Monod1}
	N.~Monod, \emph{Continuous bounded cohomology of locally compact groups}, Lecture Notes in Mathematics, No. 1758, Springer-Verlag, Berlin, 2001.
	\vspace{-7pt}
	\bibitem{Monod2}
	N.~Monod, \emph{An invitation to bounded cohomology}, International {C}ongress of {M}athematicians. {V}ol. {II}, Eur. Math. Soc., Z\"urich (2006), 1183--1211.
	\vspace{-7pt}
	\bibitem{Neeman}
	A.~Neeman, \emph{The derived category of an exact category}, J. Algebra
	\textbf{135} (1990), no.~2, 388--394.
	\vspace{-7pt}
	\bibitem{Palamodov1968}
	V.~P. Palamodov, \emph{The projective limit functor in the category of linear
		topological spaces}, Mat. Sb. (N.S.) \textbf{75 (117)} (1968), no.~4,
	567--603, English translation: Math. USSR, Sb., \textbf{4} (1968), 529--559.
	\vspace{-7pt}
	\bibitem{Palamodov1971}
	\bysame, \emph{Homological methods in the theory of locally convex spaces},
	Usp. Mat. Nauk \textbf{26} (1971), no.~1 (157), 3--65, English translation:
	Russ. Math. Surv., \textbf{26} (1971), no.~1, 1--64.
	\vspace{-7pt}
	\bibitem{BPC}
	P.~{P\'{e}rez Carreras} and J.~Bonet, \emph{Barrelled {L}ocally {C}onvex
		{S}paces}, North--Holland Mathematics Studies \textbf{113}, 1987.
	\vspace{-7pt}
	\bibitem{Prosmans2000}
	F.~Prosmans, \emph{Derived categories for functional analysis}, Publ. Res.
	Inst. Math. Sci. \textbf{36} (2000), no.~5-6, 19--83.
	\vspace{-7pt}
	\bibitem{Quillen1973}
	D.~Quillen, \emph{Higher algebraic {$K$}-theory. {I}}, Algebraic {$K$}-theory,
	{I}: {H}igher {$K$}-theories ({P}roc. {C}onf., {B}attelle {M}emorial {I}nst.,
	{S}eattle, {W}ash., 1972), Springer, Berlin, 1973, pp.~85--147. Lecture Notes
	in Math., Vol. 341.
	\vspace{-7pt}
	\bibitem{RichmanWalker}
	F.~Richman and E.~A. Walker, \emph{Ext in pre-{A}belian categories}, Pacific J.
	Math. \textbf{71} (1977), no.~2, 521--535.
	\vspace{-7pt}
	\bibitem{Rump2001}
	W.~Rump, \emph{Almost abelian categories}, Cahiers Topologie G\'eom.
	Diff\'erentielle Cat\'eg. \textbf{42} (2001), no.~3, 163--225.
	\vspace{-7pt}
	\bibitem{Rump2007}
	\bysame, \emph{A counterexample to {R}aikov's conjecture}, Bull. Lond. Math.
	Soc. \textbf{40} (2008), no.~6, 985--994.
	\vspace{-7pt}
	\bibitem{Schneiders1999}
	J.-P. Schneiders, \emph{Quasi-abelian categories and sheaves}, M\'em. Soc.
	Math. Fr. (N.S.) (1999), {\bf 76}, vi+134p..
	\vspace{-7pt}
	\bibitem{SiegThesis}
	D.~Sieg, \emph{A {H}omological {A}pproach to the {S}plitting {T}heory of
		{PLS}-spaces}, Dissertation, Universit\"at Trier, 2010.
	\vspace{-7pt}
	\bibitem{SiegWegner2009}
	D.~Sieg and S.-A. Wegner, \emph{Maximal exact structures on additive
		categories}, to appear in Math.~Nachr., 2009.
	\vspace{-7pt}
	\bibitem{Wengenroth}
	J.~Wengenroth, \emph{Derived {F}unctors in {F}unctional {A}nalysis}, Lecture
	{N}otes in {M}athematics \textbf{1810}, Springer, Berlin, 2003.
	\vspace{-7pt}
	\bibitem{Yak}
	A.~V. Yakovlev, \emph{Homological algebra in pre-abelian categories}, Zap.
	Nauchn. Sem. Leningrad. Otdel. Mat. Inst. Steklov. (LOMI) \textbf{94} (1979),
	131--141, English translation: J. Soviet Math., \textbf{19} (1982), no.~1,
	1060--1067.
	
	
	
\end{thebibliography}

{\bf \textsf{Authors' Addresses:}}
\smallskip
\\{\small Ya.~Kopylov, Sobolev Institute of Mathematics, Pr.~Akademika Koptyuga 4, 630090 Novosibirsk, RUSSIA, e-mail: yakop@math.nsc.ru.
	\smallskip
	\\S.-A.~Wegner, Fachbereich C, Bergische Universit\"at Wuppertal, Gau\ss{}stra\ss{}e 20, D-42097 Wuppertal, GERMANY, e-mail: wegner@math.uni-wuppertal.de.}





\end{document}

************************************************************************

\smallskip

In what follows, by a space we mean a locally convex space whose topology is not necessarily Hausdorff. The category formed by all these spaces with linear and continuous maps as morphisms is denoted by LCS. For notation concerning the theory of locally convex spaces we refer to \cite{FloretWloka, Jarchow, BPC}. The full subcategory of Hausdorff spaces is denoted by HD-LCS. Our first example is the category HD-BOR of Hausdorff locally convex spaces that are bornological (cf.~\cite[Example 4.2]{SiegWegner2009}), i.e.~a linear map is continuous whenever it is bounded on bounded sets. Every metrizable space is bornological and this class is stable under the formation of quotients and locally convex inductive limits. Let BOR be the full subcategory of LCS consisting of bornological spaces (cf.~\cite[\textsection{}\,23, 1.5 and \textsection{}\,11, 2.]{FloretWloka}). BOR is additive and by \cite[\textsection{}\,23, 2.9]{FloretWloka} for a morphism $f\colon E\rightarrow F$ in BOR the cokernel is the space $F/f(E)$, i.e.~the cokernels are the same as in LCS. The kernel of $f$ is the space $f^{-1}(0)^{\text{BOR}}$, i.e.~the linear space $f^{-1}(0)$ endowed with the bornological topology associated with the topology induced by $E$, see \cite[\textsection{}\,11, 2.2]{FloretWloka}. In \cite[Example 4.1]{SiegWegner2009}, it was pointed out that BOR is an example of a category which is semi-abelian but not quasi-abelian; the first property is easy to check in view of the above, whereas the absence of quasi-abelianity follows from an example due to Bonet, Dierolf \cite{BonetDierolf2005}.

\smallskip

Let us now start with an example of a category which is left semi-abelian but not right semi-abelian. The following relies on \cite[Example 4.2]{SiegWegner2009}.


\begin{ex}\label{HD-BOR} By HD-BOR we denote the full subcategory of HD-LCS consisting of bornological spaces. Clearly, HD-BOR is a full subcategory of BOR. From \cite[Example 4.1]{SiegWegner2009} it follows that the kernels in HD-BOR are those of BOR. This is not true for cokernels because a quotient space is Hausdorff if and only if the corresponding subspace is closed. Let $f\colon E\rightarrow F$ be a morphism in HD-BOR. Then the cokernel of $f$ is the space $F/\overline{f(E)}$ endowed with the quotient topology. Hence, the coimage of $f$ is the space $E/f^{-1}(0)$ endowed with the quotient topology and the image of $f$ is $\overline{f(E)}^{\scriptscriptstyle\text{BOR}}$, the closure of $f(E)$ in $F$ but endowed with the finer associated bornological topology (it is thus not clear whether $f(E)$ is still dense in this space).
\smallskip
\\Using an example due to Grothendieck \cite{Grothendieck1954} (cf.~Bonet, P\'{e}rez Carreras \cite[8.6.12]{BPC}) it was shown in \cite[Example 4.2]{SiegWegner2009} that there is a morphism $f\colon E\rightarrow F$ in HD-BOR such that $\bar{f}\colon \coim f\rightarrow \im f$ is not an epimorphism, whence HD-BOR is not right semi-abelian. On the other hand, for each morphism $f$ the induced morphism $\bar{f}$ is easily seen to be injective whence $\bar{f}$ is a monomorphism and therefore HD-BOR is left semi-abelian.
\end{ex}
\vspace{-10pt}
Let us note that the forementioned example of Grothendieck arises in the framework of so-called well-located and limit subspaces. These properties of subspaces of locally convex inductive limits arise naturally if one deals with the question whether partial differential or convolution operators on some space are surjective. We refer to Floret \cite{Floret} for precise definitions, explanations and further references.
\smallskip
\\The next example gives a category which is right semi-abelian but not left semi-abelian; our methods strongly rely on Sieg \cite{SiegThesis}.

\begin{ex}\label{HD-COM} Let HD-COM be the full subcategory of HD-LCS consisting of the complete locally convex spaces. Since finite products of complete spaces are again complete, HD-COM is additive. For a morphism $f\colon E\rightarrow F$ in HD-COM, the kernel is the space $f^{-1}(0)$ endowed with the topology induced by $E$; since we are dealing with Hausdorff spaces, $f^{-1}(0)$ is closed in $E$ and hence complete. The cokernel of $f$ is the Hausdorff completion $C(F/f(E))$, i.e.~the completion of $F/\bar{f(E)}$. Therefore, the cokernel of the kernel of $f$ is the space $C(E/f^{-1}(0))$ and the kernel of the cokernel is the space $\bar{f(E)}$ endowed with the topology induced by $F$. It is easy to see that the map $\bar{f}\colon C(E/f^{-1}(0))\rightarrow\bar{f(E)}$ has dense range and is thus an epimorphism for each morphism $f$ in HD-COM. Therefore, HD-COM is right semi-abelian. On the other hand, one can copy the proof of Sieg \cite[Proposition 3.1.6]{SiegThesis} verbatim to construct a morphism $f$ such that $\bar{f}$ is not a monomorphism. This shows that HD-COM is not left semi-abelian.
\end{ex}

\begin{ex}
By $\mathcal{C}$ we denote 
\end{ex}
As announced at the beginning of this section, we finish by reclassifying the two examples in the context of (left and right) quasi-abelian categories. For these notions, the reader is referred to \cite{Schneiders1999} and \cite{Rump2001}; note that Rump uses the term \textquotedblleft{}almost abelian\textquotedblright{} instead of \textquotedblleft{}quasi-abelian\textquotedblright{}. Using \cite[Corollary 1.2 and Proposition 1.3]{Rump2001}, we obtain a whole zoo of implications among the notions discussed so far, see Fig.~1.

It is well known that there are categories which are quasi-abelian but not abelian: For instance the category LCS, the category HD-LCS, the category of Banach spaces or the category of Fr\'{e}chet spaces are quasi-abelian but not abelian.  We refer to Prosmans \cite{Prosmans2000} for more details.

\smallskip

The question if the notions of quasi-abelian and semi-abelian categories coincide in general is known as Raikov's Conjecture and was solved in the negative by the counterexample in \cite{BonetDierolf2005} which shows that BOR is semi-abelian but not quasi-abelian (cf.~our remarks before Example \ref{HD-BOR} and \cite[Example 4.1]{SiegWegner2009}). Rump \cite{Rump2007} gave another counterexample in the context of representation theory. In fact, the category BOR also demonstrates that the implications \textquotedblleft{}left quasi-abelian $\Rightarrow$ left semi-abelian\textquotedblright{} and  \textquotedblleft{}right quasi-abelian $\Rightarrow$ right semi-abelian\textquotedblright{} can in general not be reverted: the counterexample of Bonet, Dierolf \cite{BonetDierolf2005} in fact shows that BOR is not left quasi-abelian. On the other hand, it follows directly that BOR also fails to be right quasi-abelian (assume that this is true, then \cite[Proposition 1.3]{Rump2001} would imply that BOR is even quasi-abelian). Examples \ref{HD-BOR} and \ref{HD-COM} show that all three implications starting at \textquotedblleft{}semi-abelian\textquotedblright{} and also the three implications 
ending at \textquotedblleft{}preabelian\textquotedblright{} cannot be reverted. 
\medskip
\\Let us point out that the category of PLS-spaces studied by Sieg \cite{SiegThesis} has properties similar to those of HD-COM and that the former category has important applications in analysis (see \cite{SiegThesis} for details).
\smallskip
\\To conclude, let us mention that it recently turned out that the category of LB-spaces, i.e.~countable inductive limits of Banach spaces, is an example for a category which is left quasi-abelian but not right quasi-abelian (and thus necessarily not semi-abelian). The latter was discovered by Dierolf \cite{BD} together with the second-named author.

\begin{ack}{\small Both authors would like to thank the referee for several useful comments which helped to improve this article. The first-named author was partially supported by the Russian Foundation for Basic Research (Grant 09-01-00142-a), the State Maintenance Program for the Leading Scientific Schools and Junior Scientists of the Russian Federation (NSh-6613.2010.1), and the Integration Project \textquotedblleft{}Quasiconformal Analysis and Geometric Aspects of Operator Theory\textquotedblright{} of the Siberian Branch and the Far East Branch of the Russian Academy of Sciences.}
\end{ack}
***************************************************************8
%\begin{lem}\label{1.5}\cite[Lemma 5.]{KW2} Let $\mathcal{A}$ be %semiabelian and let
%\begin{equation}
%\begin{tikzcd}[column sep=1.1ex, row sep=1.1ex]
%\cdot \arrow[]{rr}{x}\arrow[]{dd}[swap]{s}&&\cdot \arrow[]{dd}%{t}\\
%&\text{\footnotesize\rm PO}&\\
%\cdot\arrow[]{rr}[swap]{y}&&\cdot
%\end{tikzcd}
%\end{equation}
%be a pushout diagram. Assume that $x$ is strict and $\im x$ is a %semistable kernel. Then $y$ is strict and $\im y$ is a semistable %kernel.
%\end{lem}
